%!Tex root = Linguaggi.tex

\chapter{Procedure}
L'uso delle procedure permette l'\textbf{astrazione del controllo}.

\vario{Nota}{
    Astrarre significa nascondere dettagli facendo emergere caratteristcihe comuni.
}

L'astrazione chiede di identificare le proposizioni importanti che si vuole descrivere al fine di concentrarsi che si vuole descrivere al fine di concentrarsi sulle questioni rilvanti e ignorare quelle non importanti.

\smallskip
Quando si parla di astrazione del controlllo ci sono, in realtà, due categorie di sottoprogrammi con nome nome che possiamo definire:
\begin{itemize}[label=$\diamond$, itemsep=\smallskipamount]
    \item \textbf{procedure} (no output)
    \item \textbf{funzioni} (restituisce output)
\end{itemize}

\smallskip
Partre importante del processo di astrazione consiste nel dare al raggruppamento un nome.\\[0.2ex]
Il \textbf{nome} permette di rifersi all'astrazione senza conoscerne i dettagli, una solo le caratteristiche di interesse che ne hanno determinato il raggruppamento.\\[0.2ex]
I \textbf{parametri} permettono l'astrazione di tante computazioni simili che cambiano di tante computazioni simili che cambiano solo per alcuni input a cui vengono applicate e che vengono abbreviate con un singolo nome.\\[0.2ex]
Il \textbf{corpo} è ciò che viene valutato/eleborato/eseguito ogni volta che il nome viene richiamato.

\smallskip
Una procedura viene definita in tre fasi:
\begin{enumerate}[itemsep=\smallskipamount]
    \item \textbf{specifica}: descrizione dell'interfaccia della procedura/funzione
    \begin{lstlisting}[language=C]
        double P(int x) {
            
        }
    \end{lstlisting}
    \item \textbf{definizione}: descrizione del copro della procedura/funzione
    \begin{lstlisting}[language=C]
        {
            double y;
            //corpo procedura P
            ...
            (return y);
        }
    \end{lstlisting}
    \item \textbf{uso}: chiamata della procedura/funzione per eseguire il codice
    \begin{lstlisting}[language=C]
        int n;
        double z;
        z = P(n); //chiamata della procedura
    \end{lstlisting}
\end{enumerate}

Consideriamo il seguente codice:
\begin{lstlisting}[language=C, caption=programma di esempio]
    int foo(int n) {
        int tmp = a;

        if(tmp == 0)
            return n;
        else
            return n + 1;
    }
\end{lstlisting}

In questo caso la variabile \texttt{a} è un identificatore non locale, cioè un identificatore libero per cui di deve trovare il significato nel contesto.